\begin{table}[!htbp]\small\setlength{\tabcolsep}{4pt}
\caption{Composantes de la réforme de l'impôt sur les sociétés (Md€ de 2025)}\label{tab:t4}
\begin{tabularx}{\linewidth}{@{}Xrrc@{}}
\toprule
\textbf{Composante} & \textbf{2032} & \textbf{À terme} & \textbf{Statut} \\
\midrule
Fin du taux réduit de 15 % des petites entreprises & +2,0 & +2,0 & DED \\
Fin des crédits d'impôt culture (cinéma, jeux vidéo, spectacle) & +0,4 & +0,5 & DED \\
Déduction immédiate des investissements (25 % en 2029, 50 % en 2030, 100 % en 2032) & $-$7,0 & $-$1,5 & DED \\
Revalorisation des déficits reportables au taux de l'État & $-$1,1 & $-$1,1 & HYP \\
\midrule \textbf{Solde public} & \textbf{$-$5,7} & \textbf{$-$0,1} & \\
\midrule \emph{Pour mémoire, compté en partie III : recentrage du crédit d'impôt recherche} & \emph{+5,2} & \emph{+5,2} & \\
\emph{Pour mémoire, compté en partie III : mécénat, zonages, crédit famille} & & & \\
\emph{Variante : suppression du tiers restant du crédit d'impôt recherche} & \emph{+2,1} & \emph{+2,1} & HYP \\
\bottomrule
\end{tabularx}
\par\vspace{2pt}{\raggedright\scriptsize\color{gris}Source : Insee (recettes 2025) ; DGFiP, \emph{Statistiques} n\textsuperscript{o}~51 ; Évaluation des voies et moyens (PLF 2026) ; calculs des auteurs. Rendements après une décote de 10 à 25~\% pour la réaction des entreprises. La déduction immédiate coûte davantage en 2032 qu'à terme parce que l'État rembourse encore les amortissements des investissements passés. Variante : 2,6 Md€ après recentrage, décote de 20~\%.\par}
\end{table}
